\documentclass[12pt]{article}

\usepackage[a4paper, margin=1in]{geometry}
\usepackage{xcolor}
\usepackage{tcolorbox}
\usepackage{fontawesome5} % icons
\usepackage{hyperref}
\usepackage{graphicx}
\usepackage{comment}
\usepackage[nameinlink,capitalize,noabbrev]{cleveref}
\usepackage{subcaption} % modern replacement for the deprecated subfigure package

\setlength{\parindent}{0pt}

% --- COLOR PALETTE (customize these) ---
\definecolor{maincolor}{HTML}{2A9D8F} % your teal
\definecolor{accentcolor}{HTML}{E76F51}
\definecolor{lightcolor}{HTML}{FFDEAD}

% Tell cleveref how to name your counter
\crefname{exercise}{assignment}{assignments}
\Crefname{exercise}{Assignment}{Assignments}

% --- ASSIGNMENT STYLE ---
\newcounter{exercise}

% Optional label: \exercisetitle[<label>]{<Title>}
\newcommand{\exercisetitle}[2][]{%
  \refstepcounter{exercise}% make this step referenceable/hyperlinked
  \vspace{0.5cm}%
  \colorbox{maincolor}{%
    \makebox[\dimexpr\linewidth-2\fboxsep][l]{%
      \textcolor{white}{\faPen\ Assignment \theexercise:~ #2}%
    }%
  }%
  % If an optional label was provided, set it
  \if\relax\detokenize{#1}\relax\else\label{#1}\fi
  \vspace{0.3cm}%
}


% --- IMPORTANT BOX ---
\newtcolorbox{importantbox}{
  colback=lightcolor,
  colframe=accentcolor,
  arc=4mm,
  boxrule=1pt,
  left=6pt,
  right=6pt,
  top=6pt,
  bottom=6pt
}

% --- PROMPT BOX ---
\newtcolorbox{promptbox}{
  colback=white,
  colframe=maincolor,
  arc=4mm,
  boxrule=1pt,
  left=6pt,
  right=6pt,
  top=6pt,
  bottom=6pt
}

% --- QUESTION BOX (with starred variant and cleveref names) ---
% Requires: \usepackage{hyperref} for \Form/\TextField
%           \usepackage{fontawesome5} for \faEdit

% Cleveref names for the question counter
\crefname{question}{question}{questions}
\Crefname{question}{Question}{Questions}

\makeatletter
\newcounter{question}
\newlength\questionheight
\setlength\questionheight{5cm} % default answer box height

% \question{<text>}                  -> with TextField (default height)
% \question[<height>]{<text>}        -> with TextField (custom height)
% \question*{<text>}                 -> no TextField, still numbered
\newcommand{\question}{\@ifstar{\question@text}{\question@field}}

% Common heading (increments counter and creates ref anchor)
\newcommand{\question@heading}[1]{%
  \refstepcounter{question}%
  \vspace{0.3cm}%
  \noindent\textbf{\faEdit\ Question \thequestion: #1}%\\[0.3em]%
}

% Starred form: numbered question without input field
\newcommand{\question@text}[1]{%
  \question@heading{#1}%
%  \vspace{0.3cm}%
}

% Unstarred form: numbered question with input field (optional height)
\newcommand{\question@field}[2][\questionheight]{%
  \question@heading{#2}%
 \\[0.3em] \begin{Form}
    \TextField[
      name=Q\thequestion,
      multiline=true,
      width=\linewidth,
      height=#1,
      charsize=10pt,
      bordercolor={0 0 0},
      borderwidth=1.8,
      backgroundcolor={1 1 1}
    ]{}
  \end{Form}
  \vspace{0.3cm}%
}
\makeatother
% --- SOLUTION BOX ---
%--- Definition of solution/no solution env to change between the teacher and the student versions --------%

% --- Teacher toggle ---
\newif\ifteacher
\teachertrue  % Write teacherfalse for student version and teachertrue for teachers version

% --- Solution environment ---
\newtcolorbox{solutionbox}{colback=gray!10,colframe=black!60!black,arc=1mm,boxrule=1pt}

\ifteacher
\newenvironment{solution}{%
  \begin{solutionbox}\textbf{Solution:}
}{%
  \end{solutionbox}
}
\else
  \excludecomment{solution}
\fi


\begin{document}

\begin{titlepage}

% --- Top spacing ---
\vspace*{2cm}

% --- Title page ---
\begin{center}
    {\Huge \bfseries \color{maincolor} From Style to Prompt: Creating Images with an Art Historian’s Eye \par}
    \vspace{0.5cm}
    
    {\Large \color{accentcolor} Teaching with AI (TAI) \par}

    {\Large \color{accentcolor} Teachers supplementary material \par}
    
    \vspace{1.5cm}
    
    \begin{center}
{\large
Mirte van der Eyden$^{1}$, Andrea López Incera$^{2}$ \\[0.3cm]

{\small
$^{1}$ Institut für Fachdidaktik, Universität Innsbruck \\ 
mirte.van-der-eyden@uibk.ac.at \\[0.2cm]


$^{2}$ Departament de didàctica de la matemàtica i les ciències experimentals, Universitat Autònoma de Barcelona \\ 
andrea.lopez.incera@uab.cat
}
}
\end{center}

\end{center}

\vfill

% --- Logos at the bottom ---
\begin{center}
\begin{tabular}{ccc}
\includegraphics[height=1.5cm]{Figures/Logos/universitaet-innsbruck-logo-cmyk-farbe.jpg} &
\includegraphics[height=1.5cm]{Figures/Logos/logotipuab-v1-verd.png} &
\includegraphics[height=1.5cm]{Figures/Logos/EN Co-funded by the EU_POS.jpg}
\end{tabular}
\end{center}


\vspace{0.5cm}

% Optional footer text
\begin{center}
    {\small Project funded by Erasmus+ Programme of the European Union}
\end{center}

\end{titlepage}

% --- Description of activity ---
%\vspace*{0.5\textheight}
\vspace*{5 cm}
Subject: Art styles and image prompting

Activity duration: 90 min

\vspace{0.5cm}

Steps:
\begin{enumerate}
    \item Art Styles (45 minutes)
    \begin{itemize}
  \item Motivation + Assignment 1: spot the difference (15 min)
  \item Assignment 2: mystery image (30 min)
\end{itemize}
    \item Prompting images (45 minutes)
\begin{itemize}
  \item Assignment 3: reverse engineer your mystery image (10 min)
  \item Assignment 4: swap and guess (10 min)
  \item Assignment 5: create a logo (25 min)
\end{itemize}
\end{enumerate}
Learning goals:

\begin{itemize}
  \item Art styles have common visual features. 
  \item Distinguish art styles in paintings and photographs.
  \item Using AI to learn responsibly: describe images, request verification links from museums, and cross-check facts
  \item How to prompt to obtain the image you want
  \item Create a logo
\end{itemize}

Setup: 
\begin{itemize}
  \item Students need access to an AI LLM chat bot, and an AI image generation tool. 
  \item The teacher provides and prints mystery artworks for every pair or students.
  \item The teacher prints a Style Feature Guide and Style Card for every pair of students. 
\end{itemize}

\newpage

% --- HEADER / BRANDING ---
\begin{center}
    {\Huge \color{maincolor} From style to prompt}\\[5pt]
    {\large \textit{Creating images with an art historian's eye}}\\
    \rule{\linewidth}{1pt}
\end{center}

\vspace{0.5cm}

\section*{In this teacher guide you find}
\begin{itemize}
  \item Supplementary material on Assignment 1: Spot the difference
  \item List of possible mystery artworks
\end{itemize}

You can choose to focus only on paintings, or also mix in photographs in this lesson. 
In case you work with photographs as mystery art works in Assignment 2, a follow-up assignment could be that the students try to make a photograph themselves in the same style. 

\exercisetitle[]{Spot the difference}

As a warm-up, show two paintings that students compare and describe. You can choose your own pair; here is a suggested pair on the same theme (waves/sea) that differs clearly in technique, mood, and spatial depiction.
You can do the same with two photographs, suggestions are also given below. 

\subsection*{Painting A — Katsushika Hokusai, ``The Great Wave off Kanagawa'' (c.\ 1831)}
Style: Ukiyo-e (Japanese woodblock print) \\
Museum link: The Met \url{https://www.metmuseum.org/art/collection/search/45434}

\textbf{Student-friendly style features}
\begin{itemize}
  \item Clear outlines: sharp, bold lines around shapes.
  \item Flat color areas: big patches of color; little shading.
  \item Asymmetry: off-center composition; strong diagonal sweep of the wave.
  \item Detail in patterns: foam looks like claws; boats are crisp.
  \item Space feels ``flat'': Mount Fuji small and near the surface.
\end{itemize}


\subsection*{Painting B — J.\,M.\,W. Turner, ``Snow Storm — Steam-Boat off a Harbour’s Mouth'' (1842)}
Style: Romanticism \\
Museum link: Tate \url{https://www.tate.org.uk/art/artworks/turner-snow-storm-steam-boat-off-a-harbours-mouth-n00530}

\textbf{Student-friendly style features}
\begin{itemize}
  \item Loose brushwork: soft, swirling strokes; blurred edges.
  \item Atmospheric color: misty blends of gray, brown, sea tones.
  \item Dramatic light: stormy contrasts; light peeking through.
  \item Motion and mood: spinning feel; intense, emotional.
  \item Deep space: fading forms and haze suggest depth.
\end{itemize}

\textbf{Key style differences (what students can spot)}
\begin{itemize}
  \item Line vs.\ brush: Hokusai uses clear outlines; Turner dissolves edges.
  \item Flat vs.\ deep space: Hokusai looks graphic/flat; Turner suggests deep space.
  \item Color handling: Hokusai uses clean, separated colors; Turner blends into atmosphere.
  \item Detail vs.\ atmosphere: Hokusai emphasizes crisp detail; Turner emphasizes mood and movement.
  \item Mood: Hokusai is patterned/controlled; Turner is chaotic/emotional.
\end{itemize}

\textbf{Quick discussion prompts}
\begin{itemize}
  \item Where do your eyes go first in each image? Why?
  \item How does each artist make the wave feel powerful (line vs.\ brush, flat vs.\ deep space)?
  \item Which feels more like a ``poster'' or print, and which feels like a ``painting of weather''? Explain with evidence.
\end{itemize}

\subsection*{Idea for photographs to compare}

\paragraph{Photo A — Henri Cartier-Bresson, ``Behind the Gare Saint-Lazare'' (1932)}
Why it’s great for contrast: A candid, black-and-white street photograph that captures the ``decisive moment'' as a man leaps over a puddle. Unstaged, fast, and all about timing and composition. \\
Verification link: MoMA \url{https://www.moma.org/audio/playlist/244/3135}

\textbf{Style features}
\begin{itemize}
  \item Real vs.\ staged: candid/documentary; subject didn’t pose.
  \item Motion: slight blur in foot/water; action frozen mid-jump.
  \item Composition/framing: strong shapes/reflections; leading lines; off-center focus.
  \item Light/shadow: natural, soft daylight; balanced tones.
  \item Mood/message: everyday life at a perfect instant; playful/observational.
  \item Technique: 35mm film, quick shutter; classic darkroom process.
\end{itemize}


\paragraph{Photo B — Steve McCurry, ``Afghan Girl'' (1984)}
Why it’s great for contrast: A vivid color portrait for a magazine cover—staged in the sense that the subject faces the camera directly. Slow, controlled, all about gaze and color. \\
Verification links: National Geographic \url{https://www.nationalgeographic.com/photography/article/afghan-girl-steve-mccurry} \\
Steve McCurry Studio \url{https://stevemccurry.com/galleries/afghan-girl}

\textbf{Style features}
\begin{itemize}
  \item Real vs.\ staged: documentary portrait; composed with care.
  \item Color/palette: bright, intense color; green background vs.\ red headscarf; striking green eyes.
  \item Composition/framing: tight close-up; centered focus; shallow depth of field.
  \item Light/shadow: soft light enhancing texture and eye color.
  \item Mood/message: direct, intense gaze; emotional impact; social/documentary message.
  \item Technique: 35mm color film; controlled exposure and focus.
\end{itemize}


\textbf{Quick discussion prompts}
\begin{itemize}
  \item What makes A feel ``caught in the moment,'' while B feels ``confrontational and intentional''?
  \item How do black-and-white vs.\ color change the mood and what you notice first?
  \item In A, where does your eye travel? In B, what locks your gaze?
\end{itemize}


\exercisetitle[]{Mystery art work}
20 possible ``mystery'' artworks for Assignment 2. Note: website links can change over time; searching the name and artist will help locate images when the links no longer work properly.

Mix these so that each pair of students has a different style. Most images are public domain in the EU.

\begin{itemize}
  \item Early Renaissance — Sandro Botticelli, \emph{The Birth of Venus} (c.\ 1485). Tells: idealized figures, mythological subject, linear grace, shallow space. Link: Uffizi \url{https://www.uffizi.it/en/artworks/birth-of-venus}
  \item Baroque — Rembrandt, \emph{The Night Watch} (1642). Tells: dramatic light, group portrait in action, rich textures, deep space. Link: Rijksmuseum \url{https://www.rijksmuseum.nl/en/collection/SK-C-5}
  \item Baroque — Caravaggio, \emph{Supper at Emmaus} (1601). Tells: chiaroscuro, realism, theatrical gesture. Link: National Gallery (London) \url{https://www.nationalgallery.org.uk/paintings/michelangelo-merisi-da-caravaggio-the-supper-at-emmaus}
  \item Rococo — Jean-Honoré Fragonard, \emph{The Swing} (1767). Tells: pastel palette, playful sensuality, ornate detail, garden setting. Link: Wallace Collection \url{https://www.wallacecollection.org/explore/collection/search-the-collection/les-hazards-heureux-de-lescarpolette-swing/}
  \item Neoclassicism — Jacques-Louis David, \emph{Oath of the Horatii} (1784). Tells: clear contour, moral seriousness, classical subjects, rigid composition. Link: \url{https://emuseum.toledomuseum.org/objects/55069/the-oath-of-the-horatii}
  \item Romanticism — Francisco Goya, \emph{The Third of May 1808} (1814). Tells: emotion/drama, political subject, stark contrasts. Link: Prado \url{https://www.museodelprado.es/en/the-collection/art-work/the-3rd-of-may-1808-in-madrid-or-the-executions/5e177409-2993-4240-97fb-847a02c6496c}
  \item Realism — Gustave Courbet, \emph{The Stonebreakers} (1849–50; destroyed). Tells: ordinary laborers, earthy palette, anti-idealization. Link: Smarthistory \url{https://smarthistory.org/courbet-the-stonebreakers/}
  \item Impressionism — Claude Monet, \emph{Impression, Sunrise} (1872). Tells: visible brushwork, fleeting light, harbor scene, plein air. Link: \url{https://claudemonetgiverny.fr/en/actualites/impression-sunrise-the-saga-of-a-legendary-painting/}
  \item Impressionism — Pierre-Auguste Renoir, \emph{Bal du moulin de la Galette} (1876). Tells: dappled light, social leisure, vibrant brushwork. Link: Musée d’Orsay \url{https://www.musee-orsay.fr/fr/oeuvres/bal-du-moulin-de-la-galette-497}
  \item Post-Impressionism — Georges Seurat, \emph{A Sunday on La Grande Jatte} (1884–86). Tells: pointillism, optical mixing, orderly composition. Link: Art Institute of Chicago \url{https://www.artic.edu/artworks/27992}
  \item Post-Impressionism — Vincent van Gogh, \emph{The Starry Night} (1889). Tells: swirling impasto, expressive color, night landscape. Link: MoMA \url{https://www.moma.org/collection/works/79802}
  \item Post-Impressionism / Proto-Cubism — Paul Cézanne, \emph{Mont Sainte-Victoire} (c.\ 1902–06). Tells: faceted brushstrokes, structural planes, reduced depth cues. Link: The Met \url{https://www.metmuseum.org/art/collection/search/435878}
  \item Symbolism — Gustav Klimt, \emph{The Kiss} (1907–08). Tells: gold leaf, ornamental pattern, sensual theme, flat vs.\ modeled contrast. Link: Belvedere \url{https://www.belvedere.at/museum}
  \item Expressionism — Edvard Munch, \emph{The Scream} (1893). Tells: distorted forms, expressive line/color, anxiety theme. Link: Nasjonalmuseet \url{https://www.nasjonalmuseet.no/en/collection/object/NG.M.00939}
  \item Fauvism — Henri Matisse, \emph{Woman with a Hat} (1905). Tells: wild, non-natural color; loose brushwork; painterliness. Link: SFMOMA \url{https://www.sfmoma.org/artwork/91.161/}
  \item Futurism — Umberto Boccioni, \emph{Unique Forms of Continuity in Space} (1913). Tells: dynamism, movement, mechanized body, metallic flow. Link: MoMA \url{https://www.moma.org/collection/works/81179}
  \item Suprematism — Kazimir Malevich, \emph{Black Square} (1915). Tells: geometric abstraction, reduction to fundamentals, non-objective. Link: Tate \url{https://www.tate.org.uk/art/artists/kazimir-malevich-1561/five-ways-look-malevichs-black-square}
  \item De Stijl — Piet Mondrian, \emph{Broadway Boogie Woogie} (1942–43). Tells: grid, primary colors, city rhythm mapped to geometry. Link: MoMA \url{https://www.moma.org/collection/works/78682}
  \item Abstract art — Wassily Kandinsky, \emph{Composition VIII} (1923). Tells: non-objective forms, geometric/linear motifs, color–music analogy. Link: Guggenheim \url{https://www.guggenheim.org/artwork/1924}
  \item Art Nouveau — Alphonse Mucha, \emph{Job Cigarettes Poster} (1896). Tells: flowing organic lines, decorative borders, stylized hair/floral motifs. Link: Princeton University Art Museum \url{https://artmuseum.princeton.edu/art/collections/objects/11486}
  \item Ukiyo-e — Utagawa Hiroshige, \emph{Sudden Shower over Shin-Ōhashi Bridge and Atake} (1857). Tells: woodblock print, diagonal rain lines, asymmetry, flat color areas. Link: The Met \url{https://www.metmuseum.org/art/collection/search/36461}
\end{itemize}

\section*{Mystery photographs}
A curated set of 20 famous photographs for Assignment 2. Each entry includes style marks and a link (museum/academic or major agency) for research/verification.

\begin{itemize}
  \item Henri Cartier-Bresson — \emph{Behind the Gare Saint-Lazare} (1932). Style: Street / decisive moment. Marks: candid, action mid-leap, reflections/shapes, off-center framing, B\&W. Link: MoMA \url{https://www.moma.org/collection/works/98333}
  \item Dorothea Lange — \emph{Migrant Mother} (1936). Style: Documentary/social. Marks: emotional portrait, simple background, close framing, soft natural light, B\&W. Link: MoMA \url{https://www.moma.org/collection/works/50989}
  \item Ansel Adams — \emph{Clearing Winter Storm, Yosemite National Park} (c.\ 1937, printed later). Style: Landscape/modernist. Marks: dramatic mountains/clouds, strong contrast, sharp detail, deep space, B\&W. Link: AIC \url{https://www.artic.edu/artworks/121723/clearing-winter-storm-yosemite-national-park}
  \item Man Ray — \emph{Le Violon d’Ingres} (1924). Style: Surreal/studio. Marks: posed figure, playful transformation, smooth tones, careful lighting, B\&W. Link: Getty \url{https://www.getty.edu/art/collection/object/104E4A}
  \item Robert Capa — \emph{D-Day, Omaha Beach} (1944). Style: War reportage. Marks: intense action, motion blur/grain, on-the-spot framing, B\&W, urgent mood. Link: Magnum \url{https://www.magnumphotos.com/newsroom/war/robert-capa-d-day/}
  \item Cindy Sherman — \emph{Untitled Film Still \#21} (1978). Style: Staged cinematic portrait. Marks: actress-like pose, storytelling vibe, controlled framing, B\&W, classic film reference. Link: The Met \url{https://www.metmuseum.org/art/collection/search/266784}
  \item Alfred Stieglitz — \emph{The Steerage} (1907). Style: Early modernist/documentary. Marks: strong geometry (lines/decks), layered space, everyday people, B\&W. Link: The Met \url{https://www.metmuseum.org/art/collection/search/267836}
  \item Julia Margaret Cameron — \emph{Sir John Herschel} (1867). Style: Pictorialist portrait. Marks: soft focus, gentle light, artistic pose, shallow DoF, B\&W (albumen). Link: The Met \url{https://www.metmuseum.org/art/collection/search/282064}
  \item László Moholy-Nagy — \emph{Photogram} (1926). Style: Bauhaus/experimental. Marks: camera-less image, shapes/textures on paper, high contrast, abstract design, B\&W. Link: The Met \url{https://www.metmuseum.org/de/art/collection/search/265197}
  \item Gordon Parks — \emph{American Gothic, Washington, D.C.} (1942). Style: Documentary portrait/social message. Marks: staged symbolism, centered figure, B\&W, clear message. Link: Gordon Parks Foundation \url{https://www.gordonparksfoundation.org/publications/gordon-parks/american-gothic-gordon-parks-and-ella-watson}
  \item Imogen Cunningham — \emph{Magnolia Blossom} (1925). Style: Modernist close-up. Marks: tight framing, smooth light, simple shapes, elegant curves, B\&W. Link: SFMOMA \url{https://www.moma.org/collection/works/44263}
  \item Edward Weston — \emph{Pepper No.\ 30} (1930). Style: Modernist still life. Marks: single object, dramatic light, glossy texture, sculptural curves, B\&W. Link: SFMOMA \url{https://www.moma.org/audio/playlist/172/2274}
  \item Yousuf Karsh — \emph{Winston Churchill} (1941). Style: Classic studio portrait. Marks: powerful pose/expression, controlled lighting, formal clothing, centered, B\&W. Link: The Met \url{https://www.metmuseum.org/de/art/collection/search/264827}
  \item Lewis Hine — \emph{Power House Mechanic Working on Steam Pump} (c.\ 1920). Style: Industrial documentary. Marks: worker and machine, strong angles/lines, textures, B\&W, social message. Link: Brooklyn Museum \url{https://www.brooklynmuseum.org/nl-NL/objects/112132}
  \item Diane Arbus — \emph{Identical Twins, Roselle, New Jersey} (1967). Style: Portrait with uncanny mood. Marks: centered twins, flat flash, direct gaze, subtle unease, B\&W. Link: MoMA \url{https://www.moma.org/collection/works/58927}
  \item Brassaï — \emph{Paris Night View} (e.g., Street, Montparnasse, 1933). Style: Night/urban documentary. Marks: deep shadows/bright highlights, fog/wet streets, dramatic mood, B\&W. Link: Museu Tàpies \url{https://museutapies.org/en/exposicio/brassai/}
  \item Harold Edgerton — \emph{Milk Drop Coronet} (1936). Style: High-speed photography. Marks: split-second splash, bright light, clean background, scientific feel, B\&W. Link: MIT Museum \url{https://mitmuseum.mit.edu/collections/object/HEE-NC-57001}
  \item William Eggleston — \emph{Untitled (Memphis)} (c.\ 1970). Style: Color documentary/everyday. Marks: vivid color, ordinary subject, careful framing, natural light, relaxed mood. Link: The Met \url{https://www.metmuseum.org/de/art/collection/search/302301}
  \item Sebastião Salgado — \emph{Serra Pelada (Gold Mine, Brazil)} (1986). Style: Documentary/epic scale. Marks: huge crowd patterns, strong contrast, labor story, B\&W. Link: Städel Museum \url{https://sammlung.staedelmuseum.de/de/werk/goldmine-serra-pelada-state-para}
  \item Elliott Erwitt — \emph{California Kiss} (1955). Style: Street/humorous, romantic. Marks: candid moment, playful mood, simple shapes, B\&W, clear subject. Link: Koch Gallery \url{https://kochgallery.com/exhibitions/elliott-erwitt-2017/}
\end{itemize}

\section*{Tips for recreating these photography styles with a phone}
\begin{itemize}
  \item Street/doc: shoot candid moments; use burst mode; keep distance; look for strong shapes and reflections.
  \item Portrait/classic: place subject by a window; use a plain background; center or off-center per mood; convert to B\&W.
  \item Modernist close-up: find a single object; use side light; get very close; simplify the frame.
  \item Night/urban: shoot near windows or lamps; let shadows be dark; look for wet floors or shiny surfaces.
  \item High-speed: capture water splashes with bright light; use ``Live Photo'' or burst; pick the sharpest frame.
  \item Color documentary: find everyday scenes with bold colors; avoid heavy filters; let color lead the eye.
\end{itemize}



\end{document}